% ============================================================================
\clearpage
\section*{v4.9.0-r40 Source-Only: The Two-Sided Contraction across Frames, and the Weak SU(2)}
\addcontentsline{toc}{section}{v4.9.0-r40 Source-Only: The Two-Sided Contraction across Frames, and the Weak SU(2)}
\label{sec:v490-r40-two-sided-frames}
% ============================================================================

\paragraph{Purpose.}
The two-sided notation of Definition~\ref{def:two-sided-structural-tensor-interface} puts the frame on the front indices.  The author
states (2026-10-23) that its contraction discipline, Definition~\ref{def:two-sided-admissible-contractions}, has never been used or
tested, and that it is not expected to work at once.  This layer uses the front side across frames for the first time and records
the pre-registered test H25 (note \texttt{llm/Vorab\_Notiz\_H25\_Zweiseitige\_Kontraktion\_Frames.md}, commit \texttt{d1aba48}, fixed
before computing).  It also asks whether a weak \(SU(2)\) commuting with colour exists in the triality algebra.  Earlier layers are not
rewritten.

\noindent\textbf{Source-only status.}
This v4.9.0-r40 layer is source-only.  It does not build a PDF, does not assign a new DOI,
and does not authorize a public release.  The v4.9.0 Version DOI
\href{https://doi.org/10.5281/zenodo.23534954}{10.5281/zenodo.23534954} and the Concept DOI
\href{https://doi.org/10.5281/zenodo.19210728}{10.5281/zenodo.19210728} are unchanged by it.

\subsection*{v4.9.0-r40 model}
\label{subsec:v490-r40-model}

Front side: the \(\mathbb C\) factor of the Dixon storage, with the frame maps \(R_{\pi/2}\) (linear, \(X\mapsto iX\), contragredient
\(\omega\mapsto-i\omega\)), \(T\) (antilinear, complex conjugation) and \(J_B=R_{\pi/2}\circ T\).  Back side: real structure, acting on
the index \(\alpha\).  Rule~1 is the \(\mathbb C\)-bilinear contraction \(\omega_\alpha X^\alpha\); rule~2 is \(q(X)=G_{\alpha\beta}X^\alpha X^\beta\)
with real symmetric \(G\), compared with the Hermitian form \(h(X)=G_{\alpha\beta}\bar X^\alpha X^\beta\).  On
\(W=\mathbb C\otimes\HH\otimes\OO\) the front maps are \(D_4\) and left multiplication \(L_q\) by unit quaternions, the back maps right
multiplication \(R_q\) and \(SU(3)_c=\operatorname{Stab}_{G_2}(e_7)\).

\subsection*{v4.9.0-r40 results}
\label{subsec:v490-r40-results}

\paragraph{Theoremlet v4.9.0-r40.1.}
\begin{enumerate}[leftmargin=2em]
\item[(a)] On \(W\) every front map commutes with every back map.
\item[(b)] \(\omega_\alpha X^\alpha\) is invariant under \(R_{\pi/2}\) and is complex conjugated under \(T\) and \(J_B\).
\item[(c)] \(q\mapsto-q\) under \(R_{\pi/2}\), \(q\mapsto\bar q\) under \(T\), \(q\mapsto-\bar q\) under \(J_B\); \(h\) is invariant under all eight
elements of \(D_4\).
\item[(d)] The centralizer of \(\mathfrak{su}(3)_c\) in the triality algebra \(\mathfrak g\cong\mathfrak{so}(8)\) of v4.9.0-r37 is two-dimensional
and abelian.  On \(W\), \(L_q\) and \(R_q\) give two copies of \(\mathfrak{su}(2)\) that commute with each other, with
\(\mathfrak{su}(3)_c\) and with \(D_4\).
\item[(e)] No nonzero element of \(\mathfrak g\) acts trivially on one spinor slot and nontrivially on the other.
\end{enumerate}
In a check after the evaluation, the centralizer in (d) was found to be \(\operatorname{span}(L_{e_7},R_{e_7})\) on each slot.

\paragraph{Proof sketch.}
(b), (c) by direct substitution, checked on random vectors; \(D_4\) is generated by \(R_{\pi/2}\) and \(T\).  (a), (d), (e) by matrices
and null spaces (script \texttt{llm/k81\_test\_h25.py}, deviations at most \(10^{-14}\)); (e) is the injectivity of the slot projections
(v4.9.0-r37.1(a)).  \(\square\)

\paragraph{Verdict and reading.}
H25a--e are confirmed; (a) and (e) are controls.
\begin{enumerate}[leftmargin=2em]
\item Used across frames for the first time, the rules work, but not as stated.  Rule~1 says that the classical contraction remains
unchanged; across frames this holds for \(R_{\pi/2}\) only, while the antilinear frame maps conjugate the value.  Only real readouts are
frame-invariant.  Rule~2 depends on whether the comparison is bilinear or Hermitian, which the definition does not fix.  The notation
is frame-covariant, not frame-invariant; this is a finding about the notation, which so far was used in one frame only.
\item The bilinear comparison changes sign under the transition \(z\mapsto iz\) from M to T.  Read as a squared mass, this is
\(m^2\mapsto-m^2\), the tachyon.  The Hermitian comparison does not distinguish M from T.  The choice is therefore one of content:
bilinear comparisons separate the frames, Hermitian ones do not.
\item No \(\mathfrak{su}(2)\) of the triality algebra commutes with colour; what commutes is \(\operatorname{span}(L_{e_7},R_{e_7})\), the
generators of the charge \(Q=\tfrac12-\tfrac{i}{6}L_{e_7}-\tfrac{i}{3}R_{e_7}\) of v4.9.0-r35.  The triality level carries colour and
charge, but no weak interaction.  An \(\mathfrak{su}(2)\) commuting with colour appears only in the \(\HH\) factor, twice: from the left
(front side) and from the right (back side).
\item Reading: the two \(\mathfrak{su}(2)\) of \(\HH\) sit on the two sides of the two-sided notation, and ``the weak \(SU(2)\) acts only
on the left'' would read ``the weak \(SU(2)\) is one side''.  This is not shown.  The conflict of the inventory remains: the left side
also carries the Lorentz rotation, and the right side the generation \(\mathbb Z_3\) of v4.9.0-r33.
\item Proposed frame clause for rules~1 and~2 (author decision): readouts are real or are conjugated with the antilinear frame maps,
and each comparison tensor is declared bilinear (frame-separating) or Hermitian (frame-blind).
\end{enumerate}

\subsection*{v4.9.0-r40 status in the three frames}
\label{subsec:v490-r40-frames}

\begin{center}
\small
\begin{tabular}{@{}>{\raggedright\arraybackslash}p{0.30\textwidth}>{\raggedright\arraybackslash}p{0.62\textwidth}@{}}
\toprule
result & M / T / P\\
\midrule
r40.1(b), (c) & compare M with T directly: rule~1 invariant under \(R_{\pi/2}\), conjugated under \(T\), \(J_B\); bilinear rule~2 changes sign between M and T; P not evaluated\\
r40.1(a), (d), (e) & real structural statements, the same in M, T and P\\
\bottomrule
\end{tabular}
\end{center}

\subsection*{v4.9.0-r40 register entries}
\label{subsec:v490-r40-register}

\begin{center}
\small
\begin{tabular}{@{}>{\raggedright\arraybackslash}p{0.34\textwidth}>{\raggedright\arraybackslash}p{0.58\textwidth}@{}}
\toprule
statement & status after v4.9.0-r40\\
\midrule
Definition~\ref{def:two-sided-admissible-contractions} across frames & frame-covariant, not invariant (r40.1(b), (c)); frame clause proposed\\
weak \(SU(2)\) & not in the triality algebra (r40.1(d)); two candidates in \(\HH\), left and right; overload of \(\HH\) unresolved\\
colour centralizer in triality & \(\operatorname{span}(L_{e_7},R_{e_7})\), the charge generators (check after evaluation)\\
\bottomrule
\end{tabular}
\end{center}

\subsection*{v4.9.0-r40 ledger and guard update}
\label{subsec:v490-r40-ledger-guard}

\[
R_{\rm ts}^{v490\text{-}r40}=(1,0,0)
\]
for (frame behaviour of the two-sided contraction determined; frame clause undecided; weak \(SU(2)\) not located).

The layer spends no guard:
\[
R_{\rm glob}^{(v490\text{-}r40,+)}=(0,1,1,1,1),
\qquad
\Gamma_{\rm guard}^{v490\text{-}r40}=(1,1,1,1,1),
\qquad
\beta_4=1/256\ \text{unspent}.
\]
It does not imply exact finite \(C5\) zero closure, exact all-readout terminality,
\(C6/T2\), a global photon-clock/RH theorem, or global all-readout closure.
